\documentclass[11pt]{article}
\usepackage{amsmath,amssymb,amsthm,hyperref,geometry}\geometry{margin=1in}
\newtheorem{theorem}{Theorem}
\title{Structured display interaction and its native realization}
\date{15 November 2026}\begin{document}\maketitle
\paragraph{Status.} Proposed sufficient theorem for independent review. The additional structure is encoded by a specified strict 2-monad $T$ on $\mathbf{Cat}$, in fixed nested universes. This is an internalization theorem for the already specified display-exact independent criterion. It does not supply a 2-monad for arbitrary feedback, or claim that every existing contextad supports a requested structure. In particular strict algebra presentations below are an explicit restriction, not an unstated strictification theorem.

\section{Independent structured inputs}
Use the display-exact criterion with base category $C$, display class $\mathcal D$, Cartesian occurrence monad $S$, and independently valid contentad $p:L\to C$, $d:L\to C$. The letter $T$ denotes an \emph{additional} 2-monad on categories; it is not the occurrence monad $S$.

Require $C$ and $L$ to be strict $T$-algebras. Endow arrow categories and typed pullbacks of strict algebra maps with their created strict algebra structures. This construction uses $T$ on functors/natural transformations and comma/pullback universality; it does not require $T$ to preserve these limits. Assume:
\begin{enumerate}
\item The full admitted display-arrow category $K=\operatorname{Arr}_{\mathcal D}(C)$ is closed under the algebra action on $C^{\to}$, with the inherited algebra and strict inclusion. The chosen pullback/reindexing functor on its typed domain is a strict algebra morphism, and its canonical projection transformations are algebra transformations. Context unit and display accumulation are the usual identities and composition.
\item $S:C\to C$ is a strict algebra morphism; its monad unit and multiplication are algebra transformations. It preserves the display class, pullbacks, and the Cartesian unit/multiplication squares as in the unstructured criterion.
\item $p,d$ are strict algebra morphisms. The jointly typed Cartesian and opcartesian transport functors of $L$ are strict algebra morphisms; their defining universal arrows are algebra transformations. Underlying Cartesian lifts and underlying canonical invertible BC mates are those of the independent criterion.
\item The independently supplied content unit and typed accumulation are strict algebra morphisms. Their actual extension unit/multiplication bodies and coherence transformations are algebra transformations. The occurrence comparison $\nu:d\Rightarrow Sp$ is an algebra transformation satisfying the same display and unit/multiplication equations as before.
\end{enumerate}
All typed domains here are the actual categories of matching arrows and grades, formed using the strict projection maps. These assumptions specify structures on independent context and content operations. They do not mention the derived crossing, mixed comma multiplication or relative-coend algebra structure. Preservation of the chosen pullbacks is substantive: arbitrary tensor-like operations do not automatically satisfy it.

\begin{theorem}[Derived structured crossing]
The crossing of the display-exact criterion, on its full category of native representatives, is a strict $T$-algebra morphism. Its pure-middle balance, four comparison isomorphisms and all coherence transformations are algebra transformations. The underlying equations are the same full native equations as before.
\end{theorem}
\begin{proof}
The lifted context functor sends $\pi:E\to A$ to $S\pi:SE\to SA$; it is strict by the second assumption and functoriality of the arrow-category algebra. Crossing uses its pullback along $\nu_Q:dQ\to SA$. The pullback functor is strict by the first assumption. The output content $\pi^*Q$ is strict by the third assumption.

The carrier/body comparison
\[
 \chi:dQ\times_{SA}SE\longrightarrow d(\pi^*Q)
\]
is the inverse of the canonical display comparison. That comparison is an algebra transformation: its two coordinates are the extension of the defining Cartesian arrow and $\nu$, both algebra transformations by the assumptions; strictness of the pullback algebra tests its compatibility on the two projections. Its inverse is therefore an algebra transformation as well. The general representative body is obtained by additional Cartesian pullback, $\chi$, and the opcartesian body. Every ingredient is structured, so the representative crossing is strict, including on noninvertible native maps.

The canonical BC mate is structured without freely assigning it an unrelated structure. Its defining factorization is between structured functors and universal arrows. Compare its two algebra-compatibility composites after the target Cartesian arrow; the defining mate equation and structured universal arrows identify them. Both have the same base map, so Cartesian uniqueness gives equality. The same argument treats lift-composition comparisons. The four mixed comparison maps are formed from these, the structured $S$-unit/multiplication and the independently structured grade unit/accumulation. Thus they are algebra transformations. Their ten equations already hold in $\mathbf{Cat}$ by the direct criterion; the forgetful 2-functor from strict algebra morphisms is locally faithful, so they hold as structured equations. This does not use faithfulness of $d$ on grades.
\end{proof}

\section{Realization without assuming preservation of coends}
Let $H=(e\downarrow d)$ be the actual total category of mixed arrows, with source and target in $C$. Here $e:K\to C$ is the display domain functor and is strict. The structured-comma theorem (with identity comparison cells for these strict extensions) gives an algebra on $H$ with strict source and target. The pure inclusion preserves its algebra structure by the independently structured unit bodies.

\begin{theorem}[Structured native realization]
The actual mixed operational pseudocategory, with object category $C$ and arrow category $H$, is internal to strict $T$-algebras and strict algebra morphisms. Its native arrows, pure-arrow inclusion and changing-endpoint squares are the original ones. This statement neither asserts that $T$ preserves a relative pseudocoend nor replaces a category-valued module by its set of objects.
\end{theorem}
\begin{proof}
Use the direct formula for mixed composition on $H\times_C H$. Write the two mixed arrows as $f:eP\to dQ$ and $g:eR\to dV$, with $Q$ and $R$ over the common middle object. Put $K_{Q,R}=\nu_Q^*\overline S R$ and $L_{Q,R}=\epsilon_R^*Q$. The composite context is $P\otimes f^*K_{Q,R}$; its content is $V\diamond g_!L_{Q,R}$. Its body is the composite of the context accumulation body, Cartesian projection to $eK_{Q,R}$, crossing body, opcartesian body, and actual content accumulation body. Every functor and transformation in this finite formula is structured by the independent inputs and the preceding theorem. Consequently the actual multiplication functor is strict on the created algebra of composable mixed pairs.

The unitor and associator are built from one-sided coherence, reindexing, the four crossing comparisons and pure balance. These are algebra transformations by hypothesis or the preceding theorem. Their equations hold by the already proved underlying pseudocategory theorem. Thus the full pseudocategory is structured. The relative pseudocoend is used only to identify its underlying relative-pseudomonad presentation; no exchange of that coend with $T$ is needed. In particular a native square with distinct grades survives as that very square, and strict source/target retain its two changing pure endpoints.
\end{proof}

\section{Recovery and remaining weakness}
Forgetting the additional $T$-structure recovers exactly the original displayed comma construction and its pure inclusion. Whenever an original specialization admits the stated additional structure, its already established operational and cell-level recovery remains the same after forgetting. This preserves the factory's method of recovery; it does not claim every comonad, monad, Para action, span system, mixed opwreath or presheaf extension has acquired this extra structure automatically. The theorem is a sufficient structured display-family theorem within that umbrella.

The opposite-variance comma lemma permits weak extension comparisons, but the present derived-crossing theorem assumes strict primitive presentations. Moving to general colax/lax dependent transport requires a new compatibility proof, not a silent application of this result. Likewise no trace, feedback or iteration operation follows merely from calling $T$ an algebraic-structure monad. One must specify such an actual doctrine and verify the independent inputs, or choose a richer ambient when the desired structure is not captured by this 2-monadic form.

\paragraph{Relationship to existing work.} The original programme already builds on Capucci--Myers contextads and Lawler's category-valued relative setting. The structured-contextad principle in Capucci--Myers, \emph{Contextful Arrows via Wreath Products}, arXiv:2410.21889, provides the relevant one-sided architectural template: its structured-paradise and contextful-realization results already occur as Theorems 5.10 and 5.16 in v1, numbered 6.10 and 6.16 in v2. They are not a new v2 discovery. The assertion here is the stated mixed display criterion with both independent sides, its constructed crossing and its literal native comma; no claim is made that the cited theorem automatically proves these mixed hypotheses, or that this is a novel general principle of internalization.
\end{document}
